% Original STEM essay starter using essaykit (LPPL-1.3c, CTAN: essaykit).
% The prose and structure here are original catalogue content; essaykit is
% used only as the installed formatting package.
\documentclass[11pt,a4paper]{article}
\usepackage[margin=2.4cm]{geometry}
\usepackage{microtype}
\usepackage{essaykit}
\usepackage{amsmath}
\title{A Short STEM Essay}
\author{Student Name \\ Course or Institution}
\date{\today}

\begin{document}
\maketitle
\maintitle[titlecolor=mybluecolour]{Why small experiments make better engineering decisions}

\begin{abstract}
  A compact essay starter with a clear claim, evidence boxes, quotations and
  a source list. Replace the sample topic and sources with your own argument.
\end{abstract}

\section{Claim and context}
Engineering decisions are often made with incomplete information. A small
experiment narrows the uncertainty before a team commits to an expensive
design. This essay argues that the best experiment is not the most realistic
one: it is the smallest test that can distinguish between the competing
explanations for a failure or opportunity.

\rbox[title=Working claim,color=softblue]{
  Prefer experiments that change one decision, measure one observable outcome,
  and can be repeated by somebody who did not design the test.
}

\section{A practical method}
Start by writing down two hypotheses, $H_0$ and $H_1$, and the observation that
would make them disagree. If a sensor is suspected of causing latency, for
example, compare the same request path with and without that sensor while
holding the payload and load constant. The result is more useful than a large
benchmark that changes several variables at once.

\begin{enumerate}
  \item State the decision the experiment is meant to inform.
  \item Choose a measurable outcome and a stopping rule.
  \item Record the setup, including the conditions that could confound it.
  \item Run a baseline before changing the proposed mechanism.
  \item Publish the result, including an unsuccessful or inconclusive result.
\end{enumerate}

\quoter{The important thing is not to stop questioning.}{Albert Einstein}

\section{Limits and conclusion}
Small experiments do not eliminate uncertainty; they make its remaining shape
visible. A result from a quiet test environment may not predict production
behaviour, and a controlled test can miss interactions between components.
Those limits belong in the conclusion rather than in a footnote. The useful
outcome is a better next decision, not a claim that the system is proven safe.

\rrule
\sources{\href{https://www.nist.gov/}{NIST Engineering Statistics Handbook},
  \href{https://www.iso.org/standard/62085.html}{ISO 9001:2015 quality principles},
  Course notes or assigned reading}
\end{document}
